% Pursuing Stacks (cotes 134-2 to 134-8): archive cote 134-7, site page 30 (PDF page 31), typescript page 472.
% typescript. Draft compared: ps6.tex:1393-1452.
% AI-transcribed by Claude Opus 5.5 (claude-opus-5-5, adaptive thinking, effort high), Message Batches API, batch ps3 (msgbatch_01XdFmuVVBAcgc7EFxdtCqsj), 2026-11-12; not yet checked by eye.
\scanpage{134-7/30}
\mspage{472}

close as one may w\xout{h}ish to k-modules,  is to demand that moreover G is
isomorphic to the object of U(k) defined by a k-module $M_1$ ,  the group
law moreover coming from the addition law in M.  This condition is stron-%
ger still than merely demanding that G be commutative,  even when $M_1$ is
free of rank one,  \deleted{as} \inserted{because} one knows that over a non-perfect field k  there
may be "forms" of the additive group $\underline{G}_a$ which are not isomorphic to $\underline{G}_a$ ;
presumably there should be similar examples over $\underline{Z}$ to,  with rank larger
than one however.

We feel however that the case when  $G = \hand{\pi}_1(X_*)$  is a non-linear
or even a non abelian group object of U(k),  is still worth of interest.
The \xout{xxxx}\inserted[typed]{first} test it would seem,  to check if we do have a good notion indeed,
is to see if it does satisfy to the ``linearization theorem'' in case $k=\underline{Z}$ ,
i.e. the maps (w) on  page 454 are quasi-isomorphisms.  Another key test,
which now makes sense for arbitrary k,  is whether for a smooth Kan com\xout{p}-%
plex in \xout{M}$\hand{\mathcal{M}_0}(k)$  (i.e. satisfying the extra assumption involving G),
$X_*$  is homotopic to a bundle over K(G,1),  with a 1-connected fiber or
more specifically,  a fiber $Y_*$  satisfying  $Y_0=Y_1=e$ .  Among other featu-%
res to expect is a natural operation of the group object G on the k-modu-%
les $\hand{\pi}_i(X_*)$,  as well as $H_i(X_*)$ .  If however we wish,  for $k=\underline{Z}$ ,  to
use models in $\hand{\mathcal{M}_0}(\underline{Z})$  for describing possibly homotopy types with nilpo-%
tent $\hand{\pi}_1$ say, \inserted{and moreover unipotent action on the $\pi_i$} and devise a corresponding equivalence between suitable
homotopy categories,  we should first investigate the question of the
relationship between nilpotent discrete groups,  and group objects of
U(k)  -  a question already touched upon earlier in our reflection on
linearization (see end of section 94),, and of separate interest.

\deleted{\illegible{1}}

\hand{114.}  \xout{Afterxthxx}\inserted[typed]{During these} four\xout{th} days of reflection on schematization of homotopy
types,  \xout{finally} a relatively coherent picture \over{w}{h}as gradually \inserted{been} emerging from
darkness.  How far this image reflects substantial reality,  not just
daydreaming,  I would be at a loss to tell now.  Maybe some substantial
corrections will have to be made still,  besides getting in other ideas
for a more complete picture  -  I would be amazed at any rate if every-%
thing should turn out as just nonsense !  If it does'nt,  there is surely
a lot of work ahead to get everything straightened out and ready-to-use.
I will leave it at that I suppose,  for the time being  -  maybe just
finish this digression by a quick review of the set up\hand{,} and \inserted{of} some main
questions which have come out.

\margin{30}
